\section{その他}

\subsection{年金税制まとめ表}

\begin{InyouJobunbox}
  \begin{itemize}
    \item 合スト 公的年金・その他の年金制度・年金税制編 (2026版) 第9章 第3節(1)〜(2)
  \end{itemize}
\end{InyouJobunbox}

\begin{fillblank}[数理人試験 2024 問題6 設問3]
  \begin{itemize}
    \item 確定給付企業年金の加入者掛金は、（　　A　　）控除が認められる。
    \item 確定拠出企業年金の加入者掛金は、（　　B　　）控除が認められる。
  \end{itemize}
\end{fillblank}

\begin{answerbox}
  \begin{enumerate}[A:]
    \item 生命保険料
    \item 小規模企業共済等掛金
  \end{enumerate}
\end{answerbox}

\begin{truefalse}[年金1 2018 問題1(3)\textcircled{1}]
\textbf{下線部の正誤を判定し、誤っている場合は正しい内容に改めなさい。}
\\  

  確定給付企業年金の障害給付金は、年金で支給する場合、\underline{その全額が非課税}となる。
\end{truefalse}

\begin{answerbox}
  \textbf{正}
\end{answerbox}

\begin{truefalse}[年金1 2018 問題1(3)\textcircled{2}]
\textbf{下線部の正誤を判定し、誤っている場合は正しい内容に改めなさい。}
\\  

中小企業退職金共済、特定業種退職金共済、特定退職金共済、小規模企業共済のうち、掛
金を事業上の損金または必要経費に算入できない制度は\underline{特定業種退職金共済および特定退職
金共済}である。
\end{truefalse}

\begin{answerbox}
    \textbf{誤}

  小規模企業共済
\end{answerbox}


\begin{truefalse}[年金1 2018 問題1(3)\textcircled{6}]
\textbf{下線部の正誤を判定し、誤っている場合は正しい内容に改めなさい。}
\\  

退職給付信託の年金資産から稼得された、配当金（益金不算入限度超過額）や利息収入等の実際
運用収益は\underline{税務上事業主の益金および会計上の利益となる}。
\end{truefalse}

\begin{answerbox}
  \textbf{誤}

  税務上事業主の益金にはなるが、会計上の利益とはならない。
\end{answerbox}

\begin{truefalse}[年金1 2025 問題1(6)エ]
  以下の(1)から(4)のうち、正しいものをすべて挙げなさい。

【選択肢】
\begin{enumerate}[(1)]
  \item 国民年金基金に拠出した掛金は社会保険料控除の適用を受ける。
  \item 個人型確定拠出年金に拠出した掛金は一般生命保険料控除の適用を受ける。
  \item 確定拠出年金の障害給付金（年金）は雑所得として課税されるが、公的年金等控除の適用を
受ける。
  \item 厚生年金保険法に基づく遺族厚生年金は全額非課税となる。
\end{enumerate}

\end{truefalse}

\begin{answerbox}
    (1)、(4)
\end{answerbox}

\begin{explanationbox}
  \begin{enumerate}[(1)]
  \item 正しい
  \item 個人型確定拠出年金に拠出した掛金は「小規模企業共済等掛金控除」の適用を受ける。
  \item 確定拠出年金の障害給付金（年金）は「全額非課税」である。
  \item 正しい
\end{enumerate}
\end{explanationbox}


\begin{truefalse}[年金1 2019 問題2(1)\textcircled{1}]
\textbf{以下の文章の正誤を判定し、誤っている場合は正しい内容に改めなさい。}
\\  

加入者拠出のある企業型年金において、一時金として支給される老齢給付金は、加入者拠出
分を除いて一時所得として課税される。
\end{truefalse}

\begin{answerbox}
  \textbf{誤}

加入者拠出分も含めて退職所得として課税される。
\end{answerbox}

\begin{truefalse}[年金1 2019 問題2(1)\textcircled{2}]
\textbf{以下の文章の正誤を判定し、誤っている場合は正しい内容に改めなさい。}
\\  

確定給付企業年金の遺族給付金は、年金、一時金ともに所得税非課税となる。
\end{truefalse}

\begin{answerbox}
  \textbf{正}
\end{answerbox}

\begin{explanationbox}
  注意：相続税は課税される
\end{explanationbox}



\begin{truefalse}[年金1 2019 問題2(1)\textcircled{3}]
\textbf{以下の文章の正誤を判定し、誤っている場合は正しい内容に改めなさい。}
\\  

国民年金基金の遺族一時金および個人型年金の死亡一時金は、いずれも相続税の課税対象外
である。
\end{truefalse}


\begin{answerbox}
  \textbf{誤}

  個人型年金の死亡一時金は相続税の課税対象である。
\end{answerbox}



\begin{truefalse}[年金1 H29 問題1(6)\textcircled{3}]
\textbf{下線部の正誤を判定し、誤っている場合は正しい内容に改めなさい。}
\\  

国民年金基金の加入員の掛金は、\underline{小規模企業共済等掛金控除}の対象となる。
\end{truefalse}

\begin{answerbox}
  \textbf{誤}

  社会保険料控除
\end{answerbox}


\begin{truefalse}[年金1 H29 問題1(6)\textcircled{3}]
\textbf{下線部の正誤を判定し、誤っている場合は正しい内容に改めなさい。}
\\  

厚生年金基金、確定給付企業年金、確定拠出年金（企業型）のなかで、遺族給付について所
得税非課税であるのは、\underline{すべて}である。
\end{truefalse}

\begin{answerbox}
  \textbf{正}
\end{answerbox}

\begin{explanationbox}
  厚生年金基金の遺族給付は完全に非課税。

確定給付企業年金、確定拠出年金（企業型）は相続税は課税されるが、所得税は非課税。
\end{explanationbox}

\newpage

\subsection{退職給付信託の税制}


\begin{truefalse}[年金1 H29 問題1(6)\textcircled{3}]
\textbf{下線部の正誤を判定し、誤っている場合は正しい内容に改めなさい。}
\\  

退職給付信託の年金資産から稼得された、配当金（益金不算入限度超過額）や利息収入等の実際
運用収益は\underline{税務上事業主の益金および会計上の利益となる}。
\end{truefalse}

\begin{answerbox}
    \textbf{誤}


  税務上事業主の益金となるが、会計上は利益とはならない。
\end{answerbox}


\begin{truefalse}[年金1 H29 問題1(6)\textcircled{3}]
\textbf{以下の文章の正誤を判定し、誤っている場合は正しい内容に改めなさい。}
\\  

退職給付信託に拠出した資産は税務上、所有権が委託者（事業主）から受託者（信託銀行）
へ移転したものとみなすため、退職給付信託設定損益については、会計上の費用又は収益の認
識時に税務上の損金・益金処理する必要がある。
\end{truefalse}

\begin{answerbox}
  \textbf{誤}

  退職給付信託に拠出した資産は税務上、委託者である事業主が自ら保有しているものとみなすため、退職給付信託設定損益については、会計上の費用または収益の認識時に損金・益金処理することはできない。
\end{answerbox}

\newpage

\subsection{復興特別所得税}

\begin{fillblank}[年金1 2019 問題1(4)\textcircled{6}]
  東日本大震災からの復興を図ることを目的として創設された「復興特別所得税」は所得税額
の2.1％として計算され、2013 年1 月1 日から （　　A　　） 年12 月31 日までの間に支払
われる退職所得等については、所得税とともに「復興特別所得税」が源泉徴収される。
\end{fillblank}

\begin{answerbox}
  A: 2037
\end{answerbox}